\documentclass[10pt,a4paper]{amsart}
\usepackage[margin=19mm]{geometry}
\usepackage{amsmath,amssymb,mathtools}
\usepackage[scr=rsfs,cal=euler]{mathalfa}
\usepackage{xfrac}
\usepackage[unicode,hidelinks,bookmarks=false]{hyperref}
\newcommand{\ie}{i.e.\ }
\newcommand{\cf}{cf.\ }
\newcommand{\resp}{respectively\ }
\newcommand{\Subsection}{\subsection}
\providecommand{\nobreakdash}{\nobreak\hbox{-}}
\setlength{\emergencystretch}{2em}
\multlinegap0pt
\usepackage{xcolor}
\usepackage{amssymb}
\usepackage[shortlabels]{enumitem}
\usepackage{tikz}
\newtheorem{lemma}{Lemma}
\newcommand{\E}{\mathbb E}
\renewcommand{\P}{\mathbb P}
\newcommand{\R}{\mathbb R}
\renewcommand{\d}{{\rm d}}
\newcommand{\dist}{\mathop{\mathrm{dist}}}
\newcommand{\e}{\mathrm e}
\title{\textsc{Degenerate Diffusions on Continuum Percolation and Hamilton-Jacobi-Bellman Equations}}
\author{Rodrigo Bazaes, Alexander Mielke, Chiranjib Mukherjee}
\date{Source: arXiv:2606.09691v1 (CC BY 4.0)}
\begin{document}
\maketitle

\noindent\emph{Source and licence.} \textsc{Degenerate Diffusions on Continuum Percolation and Hamilton-Jacobi-Bellman Equations}. Version arXiv:2606.09691v1 (CC BY 4.0), distributed by arXiv under the Creative Commons Attribution 4.0 International licence, http://creativecommons.org/licenses/by/4.0/, as recorded in the arXiv licence metadata for this exact version and shown on the abstract page https://arxiv.org/abs/2606.09691v1. Full licence notice: \texttt{licenses/arxiv-2606.09691-CC-BY-4.0.txt}.\par\smallskip

\noindent\textit{Editorial scope.} Counted unit: the article's lemma lemma (source0.tex lines 2669--2694). The article's complete original proof of it is delivered with it (source0.tex lines 2696--3048, one \texttt{proof} environment of 353 lines). No mathematics is written, repaired or re-derived: the statement and the proof are the article's own lines, in the article's own notation, and no retained line is substituted or re-flowed.  No further source-local context is needed: every cross-reference of every retained line resolves inside the two delivered spans.  Nothing was omitted from any retained span, so the package carries no omission record.  No retained line contains a citation, so the wrapper carries no bibliography and no bibliography entry.  No retained line contains a figure environment, an image-inclusion command or drawing code, so no image is delivered and the package declares no asset dependency.  Editorial formatting only, in addition: an AMS \texttt{amsart} preamble that restates the article's own declarations and macros used by the retained lines, namely the environments \texttt{lemma} on separate counters, together with the standard packages those lines need.  The retained environments print in this document as Lemma 1 in order of appearance, whereas the article prints the same items with the numbering of its own full document.  The complete native source of the article as distributed in the arXiv e-print for this exact version is bundled unchanged as \texttt{sources/202419/source0.tex}.
\begin{lemma}
Let $\P=\P^\zeta$ denote the supercritical Boolean model and define, for all $\omega\in\Omega_0$,
\[
B(\omega):=\mathbb R^d\setminus \mathcal C_\infty(\omega),
\qquad
\delta(\omega):=
\dist\bigl(0,B(\omega)\bigr)
=
\dist\bigl(0,\partial\mathcal C_\infty(\omega)\bigr).
\]
Then there exist constants $c,C,\varepsilon_0>0$ such that for all
$0<\varepsilon<\varepsilon_0$,
\[
c\,\varepsilon
\le
\P_0(\delta<\varepsilon)
\le
C\,\varepsilon.
\]
Consequently,
\[
\E_0[\delta^{-\chi}]<\infty
\qquad\Longleftrightarrow\qquad
\chi\in (0,1).
\]
\end{lemma}

\begin{proof}
For $x\in\mathbb R^d$ and $\varepsilon>0$, define
\begin{equation}\label{def p}
p(x,\varepsilon)
:=
\P\Bigl(
0<
\dist(x,B(\omega))
<
\varepsilon
\Bigr).
\end{equation}
Since $B(\tau_x\omega)=\tau_x B(\omega)$,
we have $\dist(x,B(\omega))
=
\dist(0,B(\tau_x\omega))$ and 
by stationarity of $\P$, $p(x,\varepsilon)=p(0,\varepsilon)$.
Moreover,
\[
\{\delta>0\}
=
\{0\in\mathcal C_\infty\}
=
\Omega_0.
\]
Therefore,
\begin{equation}\label{rel}
\begin{aligned}
p(0,\varepsilon)
=
\P\Bigl(
0<
\dist(0,B(\omega))
<
\varepsilon
\Bigr)
&=
\P\bigl(
\Omega_0\cap\{\delta<\varepsilon\}
\bigr)
=
\P(\Omega_0)\,
\P_0(\delta<\varepsilon).
\end{aligned}
\end{equation}
Thus it suffices to prove that there exist constants $c,C>0$ such that
\begin{equation}\label{asymp}
c\varepsilon
\le
p(0,\varepsilon)
\le
C\varepsilon.
\end{equation}

Indeed, using
\[
\E[X^{-\chi}]
=
\chi\int_0^\infty
\varepsilon^{-\chi-1}
\P(X<\varepsilon)
\,\d\varepsilon
\]
for positive random variables $X$, and choosing
$X(\omega)=\delta(\omega)$,
we obtain
$$
\E_0[\delta^{-\chi}]
=
\chi
\int_0^\infty
\varepsilon^{-\chi-1}
\P_0(\delta<\varepsilon)
\,\d\varepsilon.
$$
Then with $\P_0(\delta<\varepsilon)\asymp\varepsilon$ 
near $0$, the integral converges if and only if $\chi<1$.
It therefore remains to prove \eqref{asymp}.

\medskip

\noindent{\bf Upper bound.}
If $0<
\dist(0,B(\omega))
<
\varepsilon$,
then there exists
$y\in\partial\mathcal C_\infty(\omega)$
such that $|y|<\varepsilon$. Since
\[
\mathcal C_\infty(\omega)
\subset
\mathcal C(\omega)
=
\bigcup_{x\in\omega}B_{1/2}(x),
\]
such a boundary point must lie on the boundary of at least one grain
$B_{1/2}(x)$, so
$|y-x|=\frac12$. Therefore,
$\Bigl||x|-\frac12\Bigr|
\leq
|y|
<
\varepsilon$, 
which implies
$x\in
B_{1/2+\varepsilon}(0)
\setminus
B_{1/2-\varepsilon}(0)$.

Hence
\[
\{0<\dist(0,B(\omega))<\varepsilon\}
\subset
\Bigl\{
\omega\cap
\bigl(
B_{1/2+\varepsilon}(0)
\setminus
B_{1/2-\varepsilon}(0)
\bigr)
\neq\emptyset
\Bigr\}.
\]

Since $\omega$ is a Poisson point process,
\[
\begin{aligned}
p(0,\varepsilon)
&\le
1-
\exp\Bigl(
-\zeta
\bigl(
|B_{1/2+\varepsilon}|
-
|B_{1/2-\varepsilon}|
\bigr)
\Bigr)
\\
&\le
C_1\varepsilon,
\end{aligned}
\]
because $|B_{1/2+\varepsilon}|
-
|B_{1/2-\varepsilon}|
=
O(\varepsilon)$.


\medskip

\noindent{\bf Lower bound. Step 1.}
Choose a small open patch
\[
\Gamma\subset S^{d-1}
\]
around the vector $e_1$ such that for all sufficiently small $\varepsilon$,
\begin{equation}\label{Seps}
S_\varepsilon
:=
\Bigl\{
x\in\mathbb R^d:
\frac12-\frac{\varepsilon}{3}<|x|<\frac12,
\quad
\frac{x}{|x|}\in\Gamma
\Bigr\}
\subset V,
\qquad
|S_\varepsilon|
=
c_0\varepsilon+O(\varepsilon^2),
\quad
c_0>0.
\end{equation}

Define the event
\[
L_\varepsilon
:=
\Bigl\{
\#(\omega\cap S_\varepsilon)=1,
\quad
\#\bigl(
\omega\cap(B_{2/3}(0)\setminus S_\varepsilon)
\bigr)=0
\Bigr\}.
\]
and note that by independence of Poisson counts on disjoint sets,
\begin{equation}\label{Leps}
\begin{aligned}
\P(L_\varepsilon)
&=
\P\bigl(
\#(\omega\cap S_\varepsilon)=1
\bigr)
\,
\P\bigl(
\#(
\omega\cap(B_{2/3}(0)\setminus S_\varepsilon)
)=0
\bigr)
=
\zeta |S_\varepsilon|
\e^{-\zeta |S_\varepsilon|}
\,
\e^{-\zeta |B_{2/3}(0)\setminus S_\varepsilon|}
\\
&=
\zeta |S_\varepsilon|
\e^{-\zeta |B_{2/3}(0)|}
=
c_1\varepsilon+O(\varepsilon^2)
\ge
c_2\varepsilon
\end{aligned}
\end{equation}
for all sufficiently small $\varepsilon$. On $L_\varepsilon$, let $x_j$ denote the unique point in $S_\varepsilon$, and define
\[
y_j
:=
-\frac{1-2|x_j|}{2|x_j|}\,x_j.
\qquad\mbox{so that}\quad 
y_j\in \partial B_{1/2}(x_j), \quad\mbox{and}\quad 
|y_j|
=
\frac{1-2|x_j|}{2}
<
\varepsilon.
\]
Moreover, for $\varepsilon<1/6$,
$B_{1/2}(y_j)\subset B_{2/3}(0)$.
Since there are no Poisson points in
$B_{2/3}(0)\setminus S_\varepsilon$,
and $x_j$ is the unique point in $S_\varepsilon$, it follows that no Poisson point other than $x_j$ lies within distance $1/2$ of $y_j$. Hence
$$y_j\in\partial\mathcal C(\omega).
$$
Finally, since $|x_j|<\frac12$,
we have
\[
0\in B_{1/2}(x_j).
\]

\medskip

\noindent{\bf Step 2.}
Choose an open ball
\[
U_1\subset \mathbb R^d\setminus B_{2/3}(0)
\quad\mbox{such that}\quad 
B_{1/2}(x)\cap B_{1/2}(u)\neq\emptyset
\quad
\forall
x\in S_\varepsilon,
\quad
u\in U_1.
\]
This is possible because $S_\varepsilon$ is contained in a fixed small sector.
Next choose finitely many pairwise disjoint open balls
\[
U_1,\dots,U_N
\subset
\R^d\setminus B_{2/3}(0)
\]
and a distant box $Q'$ such that for every $k=1,\dots,N-1$,
\[
|u_k-u_{k+1}|<1
\quad\mbox{for all}\quad 
u_k\in U_k,
\quad
u_{k+1}\in U_{k+1},
\quad\mbox{and such that}\quad 
B_{1/2}(u_N)\cap Q'\neq\emptyset
\quad \forall u_N\in U_N.
\]
If
\begin{equation}\label{J}
J
:=
\bigcap_{k=1}^N
\{\omega\cap U_k\neq\emptyset\}, \quad\mbox{then}\quad \P(J)
=
\prod_{k=1}^N
\bigl(
1-\e^{-\zeta |U_k|}
\bigr)
>0.
\end{equation}
since the sets $U_k$ are pairwise disjoint. On the event $L_\varepsilon\cap J$, the ball
$B_{1/2}(x_j)$ is connected by a finite chain of overlapping occupied balls to the box $Q'$.

\medskip

\noindent{\bf Step 3.}
Let $D$ be a bounded region containing the origin together with all the sets
\[
U_1,\dots,U_N.
\]
Choose the box
$Q'\subset D^c$
sufficiently large so that the event
\begin{equation}\label{HD}
\begin{aligned}
&H_D
:=
\bigg\{
Q'
\text{ is connected to infinity by an occupied path contained entirely in }D^c
\bigg\} \\
&\qquad\qquad\qquad \mbox{satisfies $\P(H_D)>0$}. 
\end{aligned}
\end{equation}
Such an event exists by supercriticality together with the local modification property of the Boolean model. We note that (i) the event $L_\varepsilon$ depends only on the configuration inside
$B_{2/3}(0)$; (ii)
the event $J$ depends only on the configuration inside
$D\setminus B_{2/3}(0)$;
and (iii) the event $H_D$ depends only on the configuration outside $D$.
Hence these events are independent, and therefore
\[
\P(L_\varepsilon\cap J\cap H_D)
=
\P(L_\varepsilon)\P(J)\P(H_D).
\]
On the event
$L_\varepsilon\cap J\cap H_D$,
the ball $B_{1/2}(x_j)$
is connected to infinity through occupied balls. Hence
$B_{1/2}(x_j)\subset\mathcal C_\infty(\omega)$.
Since $y_j\in\partial B_{1/2}(x_j)$
and no other ball covers $y_j$, it follows that
$y_j\in\partial\mathcal C_\infty(\omega)$.
Because $|y_j|<\varepsilon$,
we conclude that
\[
\delta(\omega)<\varepsilon.
\]
Therefore,
\[
L_\varepsilon\cap J\cap H_D
\subset
\{\delta<\varepsilon\},
\]
and thus, by \eqref{Leps}, \eqref{J} and \eqref{HD}, 
\[
\P(\delta<\varepsilon)
\ge
\P(L_\varepsilon)\P(J)\P(H_D)
\ge
c\varepsilon
\]
for some constant $c>0$ and all sufficiently small $\varepsilon$. This proves \eqref{asymp}.
\end{proof}

\end{document}
